\section{Metodologia}
\label{sec:metodologia}

Foi utilizada a amostra congelada da S02, gerada com a semente definida em
\texttt{config.yaml}, contendo 60 repositórios e cinco releases principais por
repositório (300 releases). Rascunhos e pré-releases foram excluídos. Os tipos de
projeto foram classificados como biblioteca/framework, aplicação/serviço, ferramenta
CLI ou outro. Também se registrou se as releases representam entregas reais.

Três avaliadores (A, B e C) realizaram a rotulagem manual e independente da
amostra, sem consultar os rótulos dos demais avaliadores nem a saída da heurística
avaliada. Para cada release, os avaliadores consultaram as notas de lançamento e,
quando necessário, os commits associados, classificando-a como corretiva quando a
correção de defeitos, regressões, vulnerabilidades ou a reversão de uma mudança
defeituosa constituía seu objetivo principal. As divergências foram resolvidas por
maioria simples e os casos sem maioria seriam discutidos em reunião de consenso.

A concordância foi medida pelo kappa de Fleiss e interpretada pelas faixas de
Landis--Koch~\cite{landis1977measurement}. O kappa foi 1,00 para tipo de projeto e
entregas reais e 0,742 para release corretiva, correspondente a concordância
substancial. Para releases corretivas, 80,7\% dos 300 casos foram unânimes e os 58
restantes foram decididos por maioria.

A heurística v1 classifica uma release como corretiva quando: (i) as tags atual e
anterior são SemVer estrito, admitindo apenas o prefixo \texttt{v}, possuem os mesmos
major e minor e apresentam patch crescente; ou (ii) algum assunto de commit contém o
token completo \texttt{fix}, \texttt{hotfix} ou \texttt{revert}, sem distinção entre
maiúsculas e minúsculas. Tags não-SemVer não acionam o primeiro critério, mas continuam
elegíveis ao critério textual. Foram recuperados 19.836 assuntos dos patches públicos
de 294 comparações; cinco releases não possuíam URL de comparação e uma coleta falhou.

Em relação ao consenso, a v1 obteve precisão 0,589, recall 0,906 e F1 0,714 para a
classe positiva, ultrapassando a meta de 0,70 sem refinamento adicional. A matriz de
confusão apresentou 145 verdadeiros positivos, 101 falsos positivos, 15 falsos
negativos e 39 verdadeiros negativos, com suporte positivo de 160 releases.

\begin{table}[ht]
  \centering
  \caption{Desempenho da heurística de release corretiva.}
  \label{tab:heuristica-corretiva}
  \begin{tabular}{lrrrrrrrr}
    \hline
    Versão & Precisão & Recall & F1 & Suporte & VP & FP & FN & VN \\
    \hline
    v1 & 0,589 & 0,906 & 0,714 & 160 & 145 & 101 & 15 & 39 \\
    \hline
  \end{tabular}
\end{table}

No CFR (b), uma release avaliada falha quando existe uma release classificada como
corretiva publicada em intervalo inclusivo de até sete dias. Releases publicadas nos
últimos sete dias da janela são censuradas e não integram o denominador. Assim,
$CFR_b = N_{releases\ com\ falha}/N_{releases\ avaliadas}$.

A rotulagem manual está sujeita à interpretação dos avaliadores, especialmente em
releases mistas ou com notas incompletas. Para reduzir essa ameaça, foram adotadas
definições comuns, avaliações independentes e um protocolo de consenso. Ainda assim,
o consenso não elimina completamente a subjetividade dos rótulos. O refinamento e a
avaliação da heurística sobre a mesma amostra também podem superestimar sua capacidade
de generalização. A ausência de seis comparações completas constitui uma limitação
adicional.
